\subsection{Infinitesimal transformations}

Let \(r \in \mathbf{N}_{\kappa}\) and let \(X\) be a manifold of class \(\mathbf{C}^{r+1}\) .

8.4.1. Let \(\tau\) be a vector functor for isomorphisms (7.6.6), of class \(C^r\) . If E is a Banach space, we denote by \(\mathbf{GL}(\mathbf{E})\) the open subset of \(\mathcal{L}(\mathbf{E};\mathbf{E})\) consisting of the automorphisms of E. We denote by \(\tau_{\mathbf{E}}'\) the linear tangent map at the identity element \(\mathrm{Id}_{\mathbf{E}}\) of the morphism \(\tau\colon\mathbf{GL}(\mathbf{E})\to\mathbf{GL}(\tau(\mathbf{E}))\) ; it is a continuous linear map from \(\mathcal{L}(\mathbf{E};\mathbf{E})\) to \(\mathcal{L}(\tau(\mathbf{E});\tau(\mathbf{E}))\) .

8.4.2. Let E be a Banach space and set \(F = \tau(E)\) . Let U be an open subset of E, \(\xi\) a vector field on U of class \(C^r\) , \(\tilde{\xi} \colon U \to E\) the corresponding map (obtained by identifying T(U) with U \times E (8.1.1)), and \(f \colon U \to F\) a map of class \(C^r\) . For \(x \in U\) we define an element \((\theta_\xi.f)(x)\) of F by the formula

\[
(\theta_ {\xi}. f) (x) = d _ {x} f (\tilde {\xi} (x)) - \tau_ {\mathrm{E}} ^ {\prime} (d _ {x} \tilde {\xi}) (f (x)).\tag{1}
\]

(Note that we have, on the one hand, \(\tilde{\xi}(x) \in \mathbf{E}\) and \(d_x f \in \mathcal{L}(\mathbf{E}; \mathbf{F})\), whence \(d_x f(\xi(x)) \in \mathbf{F}\); on the other hand, \(d_x \tilde{\xi} \in \mathcal{L}(\mathbf{E}; \mathbf{E})\), \(\tau_{\mathbf{E}}'(d_x \tilde{\xi}) \in \mathcal{L}(\mathbf{F}; \mathbf{F})\), and \(f(x) \in \mathbf{F}\), whence \(\tau_{\mathbf{E}}'(d_x \tilde{\xi})(f(x)) \in \mathbf{F}\). The function \(\theta_{\xi}.f: x \mapsto (\theta_{\xi}.f)(x)\) is of class \(\mathbf{C}^{r-1}\) in U.

8.4.3. Let F denote the vector bundle \(\tau(\mathrm{T}(\mathrm{X}))\) (7.6.2 and 7.6.6). Let \(\xi\) be a vector field on X and \(f\) a section of F, both of class \(\mathbf{C}^r\) . Let \(x\) be a point of X and \(c = (\mathrm{U}, \varphi, \mathrm{E})\) a chart of X at \(x\) . Let \(c' = (\mathrm{U}, \zeta_c, \mathrm{E})\) be the corresponding vector chart of T(X) (8.1.1) and \(\tau(c') = (\mathrm{U}, \psi_c, \tau(\mathrm{E}))\) the corresponding vector chart of F (7.6.2). Set \(x_c = \varphi(x)\) ; let \(\xi_c\) be the map from \(\varphi(\mathrm{U})\) into E such that \(\zeta_c(\xi(u)) = (u, \xi_c(\varphi(u)))\) and \(f_c\) the map from \(\varphi(\mathrm{U})\) into \(\tau(\mathrm{E})\) such that \(\psi_c(f(u)) = (u, f_c(\varphi(u)))\) for all \(u \in \mathrm{U}\) . Denote by \(a_c\) the element \((\theta_{\xi_c}.f_c)(x_c)\) of \(\tau(\mathrm{E})\) defined in 8.4.2. There exists an element \(a\) of \(\mathrm{F}_x\) and only one such that \(\psi_c(a) = (x, a_c)\) for every chart \(c\) of X at \(x\) . We denote it by \((\theta_\xi.f)(x)\) . It depends only on the germs of \(\xi\) and \(f\) at \(x\) (and even only on their first-order jet (12.1.2)).

The map \(\theta_{\xi}.f: x \mapsto (\theta_{\xi}.f)(x)\) is a section of class \(\mathbf{C}^{r-1}\) of F; it depends K-bilinearly on the pair \((\xi, f)\) .

8.4.4. We keep the previous hypotheses and notations. There exists (7.6.4) a morphism of vector bundles \(\tau'\) and one and only one from \(\mathcal{L}(\mathrm{T}(\mathrm{X});\mathrm{T}(\mathrm{X}))\) to \(\mathcal{L}(\mathrm{F};\mathrm{F})\) such that, for all \(x \in X\) , the restriction of \(\tau'\) to the fiber \(\mathcal{L}(T(X); T(X))_x = \mathcal{L}(T_x(X); T_x(X))\) is equal to \(\tau_{T(X)}'\) (8.4.1). If \(g\) is a function of class \(C^r\) on \(X\) , with values in \(K\) , one has

\[
\theta_ {g \xi}. f = g \theta_ {\xi}. f - \tau^ {\prime} (d g \otimes \xi) (f)
\]

(in this formula, one identifies \(dg \otimes \xi\) , which is a section of \(T'(X) \otimes T(X)\) , to a section of \(\mathcal{L}(T(X); T(X))\) ; its image \(\tau'(dg \otimes \xi)\) is a section of \(\mathcal{L}(F; F)\) and maps \(f\) to a section of \(F\) ).

8.4.5. Let I be an open subset of K containing 0, U an open subset of X, and \(\varphi\) a map of class \(C^r\) of \(\mathbf{I} \times \mathbf{U}\) in X. For \((t, x) \in \mathbf{I} \times \mathbf{U}\) , one denotes by \(\varepsilon_{t,x}\) the vector \((1, 0) \in \mathbf{K} \times \mathrm{T}_x(\mathbf{X}) = \mathrm{T}_{(t,x)}(\mathbf{I} \times \mathbf{U})\) . For \(t \in \mathbf{I}\) , one denotes by \(\varphi_t\) the map from U to X defined by \(\varphi_t(x) = \varphi(t, x) (x \in \mathbf{U})\) . We assume that the following three conditions are satisfied:

\begin{enumerate}
\def\labelenumi{\alph{enumi})}
\item
  \(\varphi_{0}\) is the canonical injection from U into X;
\item
  for all \(t \in \mathbf{I}\) , the map \(\varphi_t\) is an isomorphism of manifolds of class \(\mathbf{C}^r\) from \(\mathbf{U}\) onto an open subset of \(\mathbf{X}\) ;
\item
  the section \((t, x) \mapsto T_{(t, x)}(\varphi)(\varepsilon_{t,x})\) of the bundle \(\varphi^*T(X)\) is of class \(C^r\). Furthermore, we set:
\end{enumerate}

\[
d) \xi (x) = \mathrm{T} _ {(0, x)} (\varphi) (\varepsilon_ {0, x}) \quad \text {for all} \quad x \in \mathrm{U}.
\]

The map \(\xi\colon x\mapsto\xi(x)\) is a vector field of class \(\mathbf{C}^{r}\) on U. We say that it is the initial (or starting) field of \(\varphi\) .

Let \(\tau\) be a vector functor for isomorphisms, of class \(C^r\) ; set \(F = \tau(T(X))\) and let \(f\) be a section of class \(C^r\) of \(F\) on \(X\) . For \(x \in U\) and \(t \in I\) , set \(\varphi_t^* f(x) = \tau(T_x(\varphi_t)^{-1})(f(\varphi_t(x)))\) ; it is an element of \(F_x\) . The map \(t \mapsto \varphi_t^* f(x)\) of \(I\) to \(F_x\) is of class \(C^{r-1}\) for all \(x \in U\) and one has

\[
\frac {d}{d t} \left(\varphi_ {t} ^ {*} f (x)\right) _ {t = 0} = \left(\theta_ {\xi}. f\right) (x) \quad \text {for all} \quad x \in \mathrm{U}.\tag{2}
\]

8.4.6. Given a vector field \(\xi\) of class \(C^r\) on X and a point \(x_0\) of X, there exist an open interval I of K containing 0, an open set U of X containing \(x_0\) and a mapping \(\varphi: I \times U \to X\) of class \(C^r\) satisfying the conditions \(a\) , \(b)\) and \(c)\) of 8.4.5 and such that the initial field of \(\varphi\) is the restriction \(\xi|U\) of \(\xi\) to U.

8.4.7. Example. --- Let F be a Banach space and p an integer \(\geqslant0\) . If V is a Banach space, set \(\alpha_{p}(\mathbf{V}) = \operatorname{Alt}^{p}(\mathbf{V}; \mathbf{F})\) ; if u is an isomorphism of V onto a Banach space \(V'\) , denote by \(\alpha_{p}(u)\) the isomorphism of \(\alpha_{p}(\mathbf{V})\) onto \(\alpha_{p}(\mathbf{V}')\) deduced from u by transport of structure. One thus obtains a vector functor for isomorphisms, denoted \(\alpha_{p}\) (7.8.1 and Errata to no. 7.8). One may apply the preceding to it; one has

\[
\alpha_ {p} (\mathrm{T} (\mathrm{X})) = \operatorname{Alt}^{p}(\mathrm{T} (\mathrm{X}); \mathrm{F}).
\]

The sections of \(\alpha_{p}(\mathrm{T}(\mathrm{X}))\) are the differential forms of degree \(p\) on \(\mathbf{X}\) , with values in \(\mathbf{F}\) ; if \(\omega\) is such a form and if \(\xi\) is a vector field on \(\mathbf{X}\) , both of class \(\mathbf{C}^{r}\) , \(\theta_{\xi}.\omega\) is a differential form of degree \(p\) on \(\mathbf{X}\) , with values in \(\mathbf{F}\) and of class \(\mathbf{C}^{r-1}\) . One has

\[
\theta_ {\xi}. \omega = d (i (\xi) \omega) + i (\xi) d \omega\tag{3}
\]

and, when \(r \geqslant 2\) ,

\[
\theta_ {\xi}. d \omega = d (\theta_ {\xi}. \omega).\tag{4}
\]

When \(p=0\) , the functor \(\alpha_{p}\) is the constant vector functor defined by F; the bundle \(\alpha_{0}(\mathrm{T}(\mathrm{X}))\) is identified with the trivial bundle \(\mathrm{F}_{\mathrm{x}}\) , and the sections of this bundle are identified with the functions on X with values in F. If \(\xi\) is a vector field of class \(\mathbf{C}^{r}\) on X and if \(f\in\mathcal{C}^{r}(\mathrm{X};\mathrm{F})\) , one has

\begin{enumerate}
\def\labelenumi{(\arabic{enumi})}
\setcounter{enumi}{4}
\tightlist
\item
  \[
  \theta_ {\xi}. f = \langle \xi , d f \rangle = D _ {\xi} (f) \quad \text{(8.2.3)};
  \]
\end{enumerate}

it is an element of \(\mathcal{C}^{r-1}(X;F)\) .

8.4.8. Let \(\tau\) and \(\tau_{1}\) be two vector functors for isomorphisms. A morphism of vector functors \(h\colon \tau_{1}\to \tau\) defines a morphism of vector bundles, still denoted \(h\) , of \(\tau_{1}(\mathrm{T}(\mathbf{X}))\) to \(\tau (\mathrm{T}(\mathbf{X}))\) (7.6.4), and associates to a section \(f\) of \(\tau_{1}(\mathrm{T}(\mathbf{X}))\) a section \(h(f)\) of \(\tau (\mathrm{T}(\mathbf{X}))\) . If \(\xi\) is a vector field on \(\mathbf{X}\) and if \(f\) and \(\xi\) are of class \(\mathbf{C}^r\) , one has \(\theta_{\xi}.h(f) = h(\theta_{\xi}.f)\) .

More generally, let \(\tau\) , \(\tau_{1},\ldots,\tau_{d}\) be vector functors for isomorphisms and let \(h\colon(\tau_{1},\ldots,\tau_{d})\to\tau\) be a \(d\)-linear morphism (7.6.4). Let \(f_{i}\in\mathcal{S}_{\tau_{i}(\mathrm{T}(\mathrm{X}))}^{r}(\mathrm{X})\) (for \(i=1,\ldots,d\) ), and let \(\xi\) be a vector field of class \(\mathbf{C}^{r}\) on X. We have

\[
\theta_ {\xi}. h (f _ {1}, \dots , f _ {d}) = \sum_ {1 \leqslant i \leqslant d} h (f _ {1}, \dots , f _ {i - 1}, \theta_ {\xi}. f _ {i}, f _ {i + 1}, \dots , f _ {d}).\tag{6}
\]

For example, if \(f\) is a section of \(\tau(\mathrm{T}(\mathrm{X}))\) and \(g\) is a function on \(\mathrm{X}\) , with values in \(\mathrm{K}\) , both of class \(\mathbf{C}^r\) , one has

\[
\theta_ {\xi}. (g f) = (\theta_ {\xi}. g) f + g (\theta_ {\xi}. f).\tag{7}
\]

If \(\omega_{1}\) and \(\omega_{2}\) are two differential forms of class \(C^{r}\) , with values in Banach spaces \(F_{1}\) and \(F_{2}\) respectively, we have

\[
\theta_ {\xi}. (\omega_ {1} \wedge \omega_ {2}) = (\theta_ {\xi}. \omega_ {1}) \wedge \omega_ {2} + \omega_ {1} \wedge (\theta_ {\xi}. \omega_ {2}),\tag{8}
\]

the exterior product being taken with respect to a continuous bilinear map from \(F_{1} \times F_{2}\) into a Banach space F.

8.4.9. Let \(\tau\) be a vector functor for isomorphisms, contravariant (8.2.7), of class \(\mathbf{C}^r\) , \(\varphi: X \to Y\) a morphism of manifolds of class \(\mathbf{C}^{r+1}\) , \(\xi\) (resp. \(\eta\) ) a vector field of class \(\mathbf{C}^r\) on \(X\) (resp. \(Y\) ) such that \(\xi\) be \(\varphi\)-related to \(\eta\) (8.2.6). For every section \(f\) of \(\tau(T(Y))\) , of class \(\mathbf{C}^r\) , one has

\[
\varphi^ {*} (\theta_ {\eta}. f) = \theta_ {\xi}. \varphi^ {*} f.
\]

